% The count-emission mixture factor graph of Section~\ref{sec:emissionmixture},
% drawn by hand. Input by textbook.tex; a sketch and not a rendered figure, so
% it is outside the figure cache and the render cap, which govern generated
% output only.
\begin{figure}[htbp]
\centering
\begin{tikzpicture}[
    scale=1.0,
    within/.style={draw, thick, ->}
]
  \foreach \i in {1,...,4} {
    \node[fgvariable] (z\i) at (2.1*\i, 2.2) {$z_{\i}$};
    \node[fgdata] (e\i) at (2.1*\i, 1.4) {};
    \node[fgobserved] (n\i) at (2.1*\i - 0.45, 0.6) {$n_{\i}$};
    \node[fgobserved] (y\i) at (2.1*\i + 0.45, 0.6) {$y_{\i}$};
    \draw[fglink] (z\i) -- (e\i);
    \draw[fglink] (e\i) -- (n\i);
    \draw[fglink] (e\i) -- (y\i);
    \draw[within] (n\i) -- (y\i);
  }
  \node[fgnote] at (9.4, 2.2)
    {no pairwise factor};
  \node[fgnote] at (9.4, 1.4)
    {$\psi_i = w_{z_i}\, p(n_i, y_i \mid z_i)$};
  \node[fgnote] at (9.4, 0.6)
    {$y_i \mid n_i$ in the joint form};
\end{tikzpicture}
\caption{The count-emission mixture as a factor graph. It is
Figure~\ref{fig:mixture:sketch} with one change: the observation attached to
each unary factor is a \emph{pair}, a total $n_i$ and the successes $y_i$
within it. The arrow between them is the joint form's conditioning---the
beta-binomial's trial count is the drawn total---and is absent in the
independent form, where the trial count is a parameter and the two channels
factorize given the component. The mixture itself is unchanged either way:
Equation~\ref{eq:responsibilities} is still one normalization per variable,
and the M step is still the family's own. A sketch of the structure and not a
rendering of any instance.}
\label{fig:emissionmixture:sketch}
\end{figure}
